\documentclass[11pt]{article}
% Source-PDF: 其他 OTHERS/QTREE解法的一些研究 - Yang Zhe.pdf
\usepackage{oipl-vn}

\title{Một số nghiên cứu về lời giải cho SPOJ375 QTREE}
\author{Yang Zhe}
\date{Ngày 3 tháng 1 năm 2007}

\begin{document}
\maketitle

\begin{abstract}
Tuy bài này đã được các tác giả đi trước nghiên cứu \cite{zhou} và đã có một
lời giải khá thỏa đáng với độ phức tạp
$O(n\log n + q\sqrt n\log n)$, vẫn còn một số thuật toán tốt chưa được nhắc
đến. Bài viết này bắt đầu từ mô hình tổng quát của dạng bài này, tức bài toán
cây động, rồi giới thiệu ngắn gọn cách cải biến mô hình đó để áp dụng cho
QTREE.

Trước hết bài viết trình bày lời giải dùng trực tiếp \term{Link-Cut Trees} với
độ phức tạp thời gian $O((n+q)\log n)$. Vì lời giải này không khai thác đầy đủ
đặc điểm riêng của bài nên hiệu quả thực tế không tốt. Sau đó là một lời giải
khai thác tính chất của đề bài, có độ phức tạp $O(n+q\log^2 n)$ và chạy khá
tốt trong thực tế; trước khi tác giả nộp lời giải cuối cùng, chương trình này
xếp thứ hai trên SPOJ.

Kết quả đã có không ngăn được việc tiếp tục tìm kiếm. Khi thay cấu trúc dữ liệu
duy trì đường đi trong lời giải đó bằng \term{Splay Tree}, độ phức tạp giảm
thành $O((n+q)\log n)$, nhưng do hằng số của Splay Tree rất lớn nên hiệu quả
thực tế vẫn chưa lý tưởng. Sau khi so sánh và phân tích các lời giải trên,
bài viết đề xuất một cấu trúc dữ liệu tĩnh gọi là ``cây nhị phân cân bằng toàn
cục''. Dùng cấu trúc này để duy trì cả cây cho ta lời giải
$O((n+q)\log n)$, và hiệu quả thực tế cũng khá tốt, nhanh hơn một chút so với
thuật toán $O(n+q\log^2 n)$ đã nói ở trên.

Hy vọng bài viết có thể gợi mở thêm cho việc nghiên cứu các bài toán loại này.
\end{abstract}

\tableofcontents

\section{Mô tả bài toán và phân tích sơ bộ}

Cho một cây có $n$ đỉnh, $n\le 10000$. Mỗi cạnh có một trọng số. Cần duy trì
một cấu trúc dữ liệu hỗ trợ các thao tác sau:
\begin{enumerate}
  \item sửa trọng số của một cạnh;
  \item hỏi trọng số cạnh lớn nhất trên đường đi duy nhất giữa hai đỉnh.
\end{enumerate}
Tổng số thao tác là $q$.

Với thao tác thứ hai, chỉ cần tìm tổ tiên chung gần nhất của hai đỉnh được hỏi,
rồi lần lượt tìm trọng số cạnh lớn nhất trên hai đường đi từ hai đỉnh đó lên
tổ tiên chung. Hiển nhiên, tính trực tiếp sẽ quá chậm. Không khó nhận ra bài
toán yêu cầu một cấu trúc dữ liệu duy trì thông tin trọng số trên đường đi từ
một đỉnh lên gốc. Tuy nhiên, như vậy vẫn chưa đủ để giải quyết bài toán. Ta
hãy xét trước một bài toán tổng quát hơn: bài toán cây động.

\section{Bài toán cây động}

Bài toán cây động yêu cầu duy trì một rừng gồm nhiều cây có gốc, trong đó các
con của mỗi đỉnh không có thứ tự. Cấu trúc dữ liệu cần hỗ trợ tách cây, nối
cây, một số thao tác trên đường đi từ một đỉnh lên gốc của nó, và một số thao
tác trên cây con của một đỉnh. Nội dung đầy đủ có thể xem trong \cite{st-dt}.

Ở đây ta chỉ xét một phiên bản đơn giản hơn, chỉ gồm các thao tác thay đổi hình
dạng cây và các thao tác trên đường đi từ một đỉnh lên gốc. Nếu có thao tác
trên cây con thì cần dùng một cấu trúc khác gọi là \term{Euler-Tour Trees},
nằm ngoài phạm vi bài viết này. Cụ thể, ta duy trì một cấu trúc dữ liệu hỗ trợ:
\begin{itemize}
  \item \texttt{MAKE-TREE()} -- tạo một cây mới chỉ gồm một đỉnh;
  \item \texttt{CUT(v)} -- xóa cạnh giữa $v$ và cha của nó $\operatorname{parent}(v)$,
  tương đương tách cây con gốc $v$ ra;
  \item \texttt{JOIN(v,w)} -- cho $v$ trở thành con mới của $w$, trong đó $v$
  là gốc của một cây và $v,w$ nằm trong hai cây khác nhau;
  \item \texttt{FIND-ROOT(v)} -- trả về gốc của cây chứa $v$.
\end{itemize}

Khi hiểu rõ bài toán này, ta cũng dễ mở rộng cấu trúc dữ liệu để duy trì một
số thông tin trên đường đi từ mỗi đỉnh lên gốc của cây chứa nó, chẳng hạn tổng
trọng số, trọng số cạnh lớn nhất, hoặc độ dài đường đi.

\section{Link-Cut Trees}

\term{Link-Cut Trees} là cấu trúc dữ liệu do Sleator và Tarjan phát minh để
giải bài toán cây động \cite{st-dt}. Cấu trúc này thực hiện mỗi thao tác ở
trên trong thời gian khấu hao $O(\log n)$.

Phần này trước hết giới thiệu khái niệm Link-Cut Trees, sau đó là các thao tác
mà cấu trúc hỗ trợ, cuối cùng phân tích độ phức tạp.

\subsection{Định nghĩa Link-Cut Trees}

Ta nói một đỉnh được truy cập nếu vừa thực hiện thao tác \texttt{ACCESS} trên
đỉnh đó.

Nếu trong cây con của đỉnh $v$, đỉnh được truy cập sau cùng nằm trong cây con
gốc $w$, với $w$ là con của $v$, thì gọi $w$ là \term{Preferred Child} của
$v$. Nếu đỉnh được truy cập sau cùng chính là $v$, thì $v$ không có
Preferred Child. Cạnh nối một đỉnh với Preferred Child của nó được gọi là
\term{Preferred Edge}. Những đường đi tối đại được nối bởi Preferred Edge được
gọi là \term{Preferred Path}.

Như vậy, cả cây được chia thành một số Preferred Path. Mỗi Preferred Path được
duy trì bằng một cây cân bằng, dùng độ sâu của đỉnh trên đường đi làm khóa:
trong cây cân bằng đó, mọi đỉnh ở cây con trái của một đỉnh đều nằm phía trên
đỉnh ấy trên Preferred Path, còn mọi đỉnh ở cây con phải đều nằm phía dưới.
Cây cân bằng này phải hỗ trợ tách và hợp nhất; ở đây ta chọn Splay Tree. Ta
gọi cây cân bằng đó là một \term{Auxiliary Tree}.

Khi đã biết các Preferred Path mà cây $T$ được tách thành, chỉ cần biết thêm
quan hệ nối giữa các đường đi đó là biểu diễn được $T$. Ta dùng
\term{Path Parent} để ghi cha của đỉnh cao nhất trên Preferred Path tương ứng
với mỗi Auxiliary Tree. Nếu đỉnh cao nhất của Preferred Path chính là gốc, thì
Path Parent của Auxiliary Tree đó là \texttt{null}.

Link-Cut Trees biểu diễn mỗi cây $T$ trong rừng cần duy trì bằng nhiều
Auxiliary Tree, rồi nối các Auxiliary Tree đó bằng Path Parent.

\subsection{Các thao tác của Link-Cut Trees}

\subsubsection{\texorpdfstring{\texttt{Access}}{Access}}

\texttt{ACCESS} là nền tảng của mọi thao tác trong Link-Cut Trees. Giả sử gọi
\texttt{ACCESS(v)}. Khi đó đường đi từ $v$ lên gốc trở thành một Preferred Path
mới. Nếu trên đường đi có một đỉnh $u$ không phải Preferred Child của cha
$\operatorname{parent}(u)$, thì vì Preferred Child của $\operatorname{parent}(u)$
sẽ đổi thành $u$, Preferred Path vốn chứa $\operatorname{parent}(u)$ sẽ không
còn chứa phần từ $\operatorname{parent}(u)$ trở lên.

Hình 1 trong bản gốc minh họa một cây, Link-Cut Trees tương ứng, và trạng thái
trước/sau một thao tác \texttt{ACCESS(N)}. Văn bản trích xuất từ PDF không giữ
được bố cục hình, nên bản dịch này chỉ giữ nội dung thuật toán.

Cách thực hiện \texttt{ACCESS} có vẻ bất ngờ: ta cập nhật trực tiếp các
Auxiliary Tree cần thay đổi. Thoạt nhìn cách này không hiệu quả, nhưng phần
sau sẽ chứng minh độ phức tạp khấu hao của nó là $O(\log n)$.

Trước hết, vì đã truy cập $v$, Preferred Child của $v$ phải biến mất. Ta xoay
$v$ lên gốc của Auxiliary Tree chứa nó. Nếu $v$ có con phải trong Auxiliary
Tree đó, tức Preferred Child cũ của $v$, thì tách cây con phải của $v$ ra khỏi
Auxiliary Tree hiện tại, và đặt Path Parent của Auxiliary Tree mới này là
$v$.

Sau đó, nếu Preferred Path chứa $v$ chưa chứa gốc của cây thật, gọi Path Parent
của nó là $u$. Ta xoay $u$ lên gốc của Auxiliary Tree chứa $u$, thay cây con
phải của $u$ trong Auxiliary Tree đó bằng Auxiliary Tree chứa $v$, rồi đặt
Path Parent của cây con phải cũ của $u$ thành $u$. Lặp quá trình này cho đến
khi gặp Preferred Path chứa gốc.

\subsubsection{\texorpdfstring{\texttt{Find Root}}{Find Root}}

Sau \texttt{ACCESS(v)}, gốc chắc chắn là đỉnh nhỏ nhất trong Auxiliary Tree
chứa $v$. Trước hết xoay $v$ lên gốc của Auxiliary Tree đó. Sau đó bắt đầu từ
$v$ đi liên tục sang trái trong Auxiliary Tree cho đến khi không thể đi tiếp;
đỉnh đó chính là gốc cần tìm. Vì Auxiliary Tree được lưu bằng Splay Tree, ta
còn thực hiện \texttt{Splay} trên đỉnh gốc vừa tìm được.

\subsubsection{\texorpdfstring{\texttt{Cut}}{Cut}}

Trước hết truy cập $v$, rồi xoay $v$ lên gốc của Auxiliary Tree chứa nó. Sau
đó ngắt liên kết giữa $v$ và cây con trái của nó trong Auxiliary Tree.

\subsubsection{\texorpdfstring{\texttt{Join}}{Join}}

Trước hết truy cập $v$, rồi sửa Path Parent của Auxiliary Tree chứa $v$ thành
$w$, sau đó truy cập lại $v$.

\subsubsection{Mã giả của các thao tác}

\begingroup
\small
\begin{verbatim}
procedure Access(v)
    u <- v
    v <- null
    repeat
        Splay(u)
        path-parent[right-child[u]] <- u
        right-child[u] <- v
        path-parent[v] <- null
        v <- u
        u <- path-parent[u]
    until u = null
end procedure

function FindRoot(v)
    Access(v)
    Splay(v)
    while left-child[v] != null do
        v <- left-child[v]
    end while
    Splay(v)
    return v
end function

procedure Cut(v)
    Access(v)
    left-child[v] <- null
end procedure

procedure Join(v, w)
    Access(v)
    path-parent[v] <- w
    Access(v)
end procedure
\end{verbatim}
\endgroup

\subsection{Phân tích độ phức tạp của Link-Cut Trees}

Từ mã giả có thể thấy, ngoài \texttt{ACCESS}, mọi thao tác khác có độ phức tạp
khấu hao nhiều nhất $O(\log n)$. Vì vậy chỉ cần phân tích độ phức tạp của
\texttt{ACCESS}.

Trong phân tích thao tác \texttt{ACCESS}, ta chia thời gian thành hai phần.
Thứ nhất, chứng minh số lần đổi Preferred Child là $O(\log n)$ khấu hao. Thứ
hai, chứng minh tổng thời gian của mọi thao tác Splay trong một lần
\texttt{ACCESS} là $O(\log n)$ khấu hao.

\subsubsection{Phân rã đường đi nặng-nhẹ}

\term{Phân rã đường đi nặng-nhẹ} là một kỹ thuật phân tích áp dụng cho mọi cây.
Nó chia các cạnh của cây thành hai loại: cạnh nặng và cạnh nhẹ.

Ký hiệu $\operatorname{size}(i)$ là số đỉnh trong cây con gốc $i$. Gọi $u$ là
đứa con có $\operatorname{size}(u)$ lớn nhất trong các con của $v$; nếu có
nhiều đứa con như vậy thì chọn tùy ý. Cạnh $(v,u)$ được gọi là cạnh nặng, còn
các cạnh $(v,w)$ từ $v$ đến những đứa con khác $w$ được gọi là cạnh nhẹ, ngay
cả khi $\operatorname{size}(w)=\operatorname{size}(u)$. Mỗi đỉnh có con thì có
duy nhất một cạnh nặng đi ra từ nó. Từ định nghĩa này suy ra ngay:

\paragraph{Tính chất 3.1.}
Nếu $u$ là con của $v$ và $(v,u)$ là một cạnh nhẹ, thì
\[
  \operatorname{size}(u) < \operatorname{size}(v)/2.
\]

\paragraph{Chứng minh.}
Giả sử $\operatorname{size}(u)\ge \operatorname{size}(v)/2$. Vì $(v,u)$ là
cạnh nhẹ, tồn tại một cạnh nặng $(v,w)$ đi ra từ $v$. Theo định nghĩa cạnh
nặng,
\[
  \operatorname{size}(v)
  \ge 1+\operatorname{size}(w)+\operatorname{size}(u)
  \ge 2\operatorname{size}(u)+1
  \ge 1+\operatorname{size}(v),
\]
mâu thuẫn.

Gọi $\operatorname{light\mbox{-}depth}(v)$ là số cạnh nhẹ trên đường từ $v$ lên
gốc.

\paragraph{Bổ đề 3.1.}
\[
  \operatorname{light\mbox{-}depth}(v) \le \lg n.
\]

\paragraph{Chứng minh.}
Theo Tính chất 3.1, mỗi khi đi qua một cạnh nhẹ thì kích thước cây con hiện tại
giảm xuống nhiều nhất còn một nửa. Ban đầu kích thước cây con không quá $n$,
còn cây con của $v$ có ít nhất một đỉnh, nên đi qua nhiều nhất $\lg n$ cạnh
nhẹ.

\subsubsection{Chứng minh số lần đổi Preferred Child là
\texorpdfstring{$O(\log n)$}{O(log n)} khấu hao}

Để đếm số lần Preferred Child thay đổi, trước hết ta thực hiện phân rã nặng-nhẹ
trên cây $T$.

Ngoài việc Preferred Child của $v$ biến mất, mỗi lần Preferred Child thay đổi
đều sinh ra một Preferred Edge mới. Vì vậy chỉ cần đếm số Preferred Edge mới
được sinh ra.

Theo Bổ đề 3.1, trong quá trình \texttt{ACCESS} có nhiều nhất $\lg n$ Preferred
Edge nhẹ mới. Tuy nhiên, một lần truy cập có thể sinh ra nhiều Preferred Edge
nặng. Ta đếm chúng như sau. Số lần một cạnh nặng trở thành Preferred Edge nhiều
nhất bằng số lần nó từ Preferred Edge trở lại cạnh thường, cộng với tổng số
cạnh, vì cuối cùng có thể còn lại toàn bộ các cạnh nặng là Preferred Edge.
Mỗi lần một cạnh nặng từ Preferred Edge trở lại cạnh thường thì tương ứng có
một cạnh nhẹ trở thành Preferred Edge. Do đó tổng số lần cạnh nặng trở thành
Preferred Edge nhiều nhất bằng tổng số lần cạnh nhẹ trở thành Preferred Edge
cộng với tổng số cạnh. Trung bình trên mỗi thao tác \texttt{ACCESS}, số cạnh
nặng trở thành Preferred Edge là $O(\log n)$.

Tóm lại, số lần Preferred Child thay đổi là $O(\log n)$ khấu hao.

\subsubsection{Chứng minh tổng thời gian Splay là
\texorpdfstring{$O(\log n)$}{O(log n)} khấu hao}

Phân tích độ phức tạp của Splay thường dùng phương pháp thế năng.

Gán cho mỗi đỉnh $u$ trong Splay Tree một trọng số dương $w(u)$, và gọi $s(u)$
là tổng trọng số các đỉnh trong cây con của $u$. Định nghĩa hàm thế năng
\[
  \Phi=\sum_u \lg s(u).
\]
Các kết quả chuẩn sau đây được dùng; trong đó $s(u)$ là giá trị trước thao tác
và $s'(u)$ là giá trị sau thao tác.

\paragraph{Bổ đề 3.2.}
Trong thao tác \texttt{SPLAY(v)}, chi phí khấu hao của bước zig-zig hoặc
zag-zag thỏa
\[
  \widehat{\operatorname{cost}}
  \le 3\bigl(\lg s'(v)-\lg s(v)\bigr).
\]
Hằng số $3$ ở đây không thể cải thiện.

\paragraph{Bổ đề 3.3.}
Trong thao tác \texttt{SPLAY(v)}, chi phí khấu hao của bước zig-zag hoặc
zag-zig thỏa
\[
  \widehat{\operatorname{cost}}
  \le 2\bigl(\lg s'(v)-\lg s(v)\bigr)
  \le 3\bigl(\lg s'(v)-\lg s(v)\bigr).
\]

\paragraph{Bổ đề 3.4.}
Trong thao tác \texttt{SPLAY(v)}, chi phí khấu hao của bước zig hoặc zag thỏa
\[
  \widehat{\operatorname{cost}}
  \le \lg s'(v)-\lg s(v)+1
  \le 3\bigl(\lg s'(v)-\lg s(v)\bigr)+1.
\]

Chứng minh ba bổ đề này khá đơn giản; độc giả quan tâm có thể tự chứng minh
hoặc tham khảo tài liệu liên quan. Từ đó suy ra:

\paragraph{Định lý 3.5 (Access Theorem).}
Chi phí khấu hao của \texttt{SPLAY(v)} thỏa
\[
  \widehat{\operatorname{cost}}
  \le 3\bigl(\lg s'(v)-\lg s(v)\bigr)+1.
\]
Hằng số $3$ ở đây không thể cải thiện.

Trong bài toán hiện tại, vì một thao tác \texttt{ACCESS} thực hiện nhiều thao
tác \texttt{SPLAY} (khấu hao $O(\log n)$ lần), ta cần chọn một hàm trọng số
phù hợp, thay vì đơn giản đặt $w(u)=1$. Nếu xem các cạnh trong Auxiliary Tree
là cạnh thật, còn Path Parent là cạnh ảo, thì tổng số đỉnh của một cây con bằng
$1$ cộng với tổng số đỉnh của mọi cây con của nó, kể cả các cây con nối bằng
cạnh ảo. Đặt
\[
  w(u)=1+\text{tổng số đỉnh trong các cây con có Path Parent là }u,
\]
khi đó $s(u)$ chính là số đỉnh trong cây con của $u$.

Giả sử trong quá trình \texttt{ACCESS(v)}, ta lần lượt thực hiện Splay trên
$v_0=v$, $v_1=\operatorname{path\mbox{-}parent}[v_0]$, \ldots,
$v_k=\operatorname{path\mbox{-}parent}[v_{k-1}]$. Khi đó chi phí khấu hao của
thao tác \texttt{ACCESS} là
\[
\begin{aligned}
\widehat{\operatorname{cost}}
&\le \sum_{i=0}^{k} 3\bigl(\lg s'(v_i)-\lg s(v_i)\bigr)+1 \\
&\le
3\left[
\sum_{i=1}^{k}\bigl(\lg s(v_i)-\lg s(v_{i-1})\bigr)
+ \lg s(v_0)-\lg s(v_0)
\right]+k \\
&= 3\bigl(\lg s(v_k)-\lg s(v_0)\bigr)+k \\
&\le 3\lg n + \text{số lần Preferred Child thay đổi}.
\end{aligned}
\]
Ở dòng giữa, bản PDF bị hỏng chỉ số một phần; công thức trên giữ đúng ý nghĩa
telescoping của chứng minh trong bài gốc.

Vì đã chứng minh số lần Preferred Child thay đổi là $O(\log n)$ khấu hao, thao
tác \texttt{ACCESS} cũng có độ phức tạp khấu hao $O(\log n)$.

\section{Lời giải cho bài toán}

\subsection{Tổng quan các lời giải}

Với các kiến thức trên, ta có thể giải tốt bài toán. Bài viết thảo luận bốn
lời giải, như Bảng~\ref{tab:solutions}.

\begin{table}[htbp]
\centering
\begin{tabular}{cp{0.56\linewidth}c}
\hline
\textbf{Số} & \textbf{Tóm tắt lời giải} & \textbf{Độ phức tạp} \\
\hline
1 & Cây động, dùng Link-Cut Trees & $O((n+q)\log n)$ \\
2 & Phân rã nặng-nhẹ, dùng cây đoạn hoặc cây nhị phân ảo cho từng đường &
$O(n+q\log^2 n)$ \\
3 & Phân rã nặng-nhẹ, dùng Splay Tree cho từng đường & $O((n+q)\log n)$ \\
4 & Phân rã nặng-nhẹ, dùng một cây nhị phân cân bằng toàn cục cho mọi đường &
$O((n+q)\log n)$ \\
\hline
\end{tabular}
\caption{Các lời giải được trình bày trong bài}
\label{tab:solutions}
\end{table}

Các lời giải trong bài gốc đều có chương trình tham khảo. Thời gian chạy của
chúng được ghi trong Bảng~\ref{tab:runtime}.

\begin{table}[htbp]
\centering
\begin{tabular}{lp{0.58\linewidth}c}
\hline
\textbf{Lời giải} & \textbf{Chương trình tham khảo} & \textbf{Thời gian} \\
\hline
1 & \texttt{QTREEdynamic-tree link-cut-trees.pas} & 4.58 giây \\
2, cây đoạn & \texttt{QTREEheavy-light-decomposition segment-tree.pas} &
2.65 giây \\
2, cây nhị phân ảo & \texttt{QTREEheavy-light-decomposition imaginary-bst.pas} &
2.37 giây \\
3 & \texttt{QTREEheavy-light-decomposition splay-tree.pas} & 4.28 giây \\
4 & \texttt{QTREEheavy-light-decomposition global-balanced-bst.pas} &
2.06 giây \\
\hline
\end{tabular}
\caption{So sánh thời gian chạy của các chương trình}
\label{tab:runtime}
\end{table}

Cần chú ý rằng lời giải 2 và 3 tuy chỉ khác cấu trúc dữ liệu dùng để duy trì
đường đi, nhưng vì độ phức tạp thời gian khác nhau, và việc so sánh chúng đóng
vai trò quan trọng trong việc dẫn tới lời giải 4, nên được liệt kê riêng. Hai
chương trình của lời giải 2 chỉ khác cách cài đặt; thời gian chạy gần nhau và
không đủ để đánh giá cây đoạn hay cây nhị phân ảo tốt hơn.

\subsection{Lời giải cây động}

Dùng Link-Cut Trees để duy trì cây. Ta chỉ cần mở rộng Auxiliary Tree đã nói ở
trên: trong Auxiliary Tree, với mỗi đỉnh ghi trọng số cạnh nối nó với cha là
\texttt{cost}, và ghi \texttt{maxcost}, tức giá trị \texttt{cost} lớn nhất
trong cây con gốc đỉnh đó. Vì phép lấy lớn nhất có tính kết hợp, \texttt{maxcost}
có thể được duy trì trong $O(1)$ khi xoay Auxiliary Tree.

Với thao tác sửa trọng số cạnh, chỉ cần cập nhật \texttt{cost} của đỉnh tương
ứng trong Auxiliary Tree chứa nó, rồi Splay đỉnh đó để duy trì các giá trị
\texttt{maxcost} liên quan trong Auxiliary Tree. Rõ ràng thao tác này mất
$O(\log n)$.

Với truy vấn trọng số cạnh lớn nhất trên đường đi từ $u$ đến $v$, trước hết
\texttt{ACCESS(u)}. Sau đó, trong quá trình \texttt{ACCESS(v)}, khi đi tới
Auxiliary Tree chứa gốc (tức chứa $u$), đỉnh đang xét đúng là tổ tiên chung gần
nhất của $u$ và $v$; gọi đỉnh đó là $d$. Khi đó \texttt{maxcost} của cây con
phải của $d$ trong Auxiliary Tree chứa $d$ là trọng số lớn nhất trên đường từ
$u$ đến $d$, còn \texttt{maxcost} của Auxiliary Tree chứa $v$ là trọng số lớn
nhất trên đường từ $v$ đến $d$. Đáp án là giá trị lớn hơn trong hai
\texttt{maxcost} đó. Vì \texttt{Access} có độ phức tạp khấu hao $O(\log n)$,
truy vấn cũng mất $O(\log n)$.

Do thế năng ban đầu của Link-Cut Trees là $O(n\log n)$, tổng độ phức tạp của
phương pháp này là $O((n+q)\log n)$. Tuy nhiên, vì dùng Splay Tree làm cấu trúc
cho Auxiliary Tree, mà hằng số của Splay Tree rất lớn, thuật toán này không dễ
vượt giới hạn thời gian 5 giây. Kể cả tác giả của \cite{zhou} cũng không dùng
cách này để qua toàn bộ dữ liệu. Phần tiếp theo mô tả một mẹo cài đặt Splay
Tree trong chương trình của tác giả.

\subsubsection{Một mẹo cài đặt Splay Tree}

Trong Định lý 3.5, ta thấy hằng số của Splay Tree bằng ba lần thời gian của
một phép xoay đơn. Trong một phép xoay đơn, không chỉ phải sửa hình dạng cục bộ
của Splay Tree mà còn phải duy trì thuộc tính \texttt{maxcost}; viết tốt phần
này rất quan trọng với hiệu quả thực tế của Splay Tree. Đồng thời, cách viết
Splay Tree thông thường phải kiểm tra đỉnh có phải gốc không, hoặc là con nào
của cha, cũng tốn nhiều thời gian. Vì vậy tác giả tối ưu các thao tác này và
thu được cách viết \texttt{SPLAY}, \texttt{ROTATE}, \texttt{UPDATE} vừa ngắn
gọn vừa hiệu quả.

Giải thích một số biến: để tiện cài đặt \texttt{UPDATE}, đặt một đỉnh rỗng giả
\texttt{null} với \texttt{maxcost = -infinity}. Để tiện cài đặt
\texttt{ROTATE}, giả sử con trái của đỉnh $x$ lưu ở \texttt{child[x][0]}, con
phải lưu ở \texttt{child[x][1]}. Nếu $x$ không phải gốc, \texttt{parent[x]} lưu
cha của $x$, còn \texttt{nodetype[x]} ghi $x$ là con nào của cha, tức
\texttt{child[parent[x]][nodetype[x]] = x}. Nếu $x$ là gốc, \texttt{parent[x]}
lưu Path Parent của Auxiliary Tree này và \texttt{nodetype[x] = 2}.

\begingroup
\small
\begin{verbatim}
procedure Update(node)        // cap nhat maxcost cua node
    k <- cost[node]
    x <- maxcost[child[node][0]]
    y <- maxcost[child[node][1]]
    if x < y then
        x <- y
    end if
    if k < x then
        k <- x
    end if
    maxcost[node] <- k
end procedure

procedure Rotate(node, x)     // xoay node co nodetype x len vi tri cua cha
    p <- parent[node]
    y <- nodetype[p]

    tmp <- parent[p]
    parent[node] <- tmp
    nodetype[node] <- y
    if y != 2 then
        child[tmp][y] <- node
    end if

    y <- 1 - x
    tmp <- child[node][y]
    child[p][x] <- tmp
    parent[tmp] <- p
    nodetype[tmp] <- x

    parent[p] <- node
    nodetype[p] <- y
    child[node][y] <- p

    Update(p)
end procedure

procedure Splay(node)
    repeat
        a <- nodetype[node]
        if a = 2 then
            break
        end if
        p <- parent[node]
        b <- nodetype[p]

        if a = b then
            Rotate(p, a)       // buoc dau cua zig-zig hoac zag-zag
        else
            Rotate(node, a)    // buoc dau cua zig-zag, zag-zig, zig hoac zag
        end if
        if b = 2 then
            break              // p la goc cu, node da duoc xoay len goc
        end if
        Rotate(node, b)        // buoc thu hai
    until false
    Update(node)
end procedure
\end{verbatim}
\endgroup

Trong quá trình \texttt{SPLAY(node)}, cách viết thông thường sẽ gọi
\texttt{UPDATE} trên một số đỉnh quá nhiều lần, nhưng thực ra không cần thiết.
Chỉ cần trong \texttt{ROTATE} cập nhật cha của đỉnh được xoay, rồi cuối
\texttt{SPLAY} cập nhật một lần trên \texttt{node}. Vì \texttt{UPDATE} là thao
tác tốn thời gian, tối ưu này loại bỏ khoảng một nửa số lần \texttt{UPDATE}
thừa và cải thiện đáng kể hiệu năng chương trình.

\subsection{Phương pháp phân rã đường đi nặng-nhẹ}

Trong bài này, hình dạng cây không thay đổi, nên ta hoàn toàn không dùng tới
các thao tác thay đổi hình dạng cây trong bài toán cây động. Với một bài toán
đặc biệt hơn như vậy, đáng lẽ phải có một thuật toán không phức tạp hơn
Link-Cut Trees và không chậm hơn $O((n+q)\log n)$.

Link-Cut Trees dùng các khái niệm như Preferred Child để cho phép hình dạng cây
thay đổi. Liệu ta có thể bỏ các khái niệm đó không? Sau khi phân tích, ta thấy
có thể. Chỉ cần áp dụng kỹ thuật phân rã đường đi nặng-nhẹ trong phần phân
tích độ phức tạp của Link-Cut Trees là thu được lời giải nhanh cho bài này.

\subsection{Dùng cây đoạn hoặc cây nhị phân ảo để duy trì đường đi}

Trước hết, thực hiện phân rã nặng-nhẹ trên cây. Gọi một đường đi tối đại gồm
các cạnh nặng là \term{đường nặng}. Khi đó cây được chia thành nhiều đường
nặng. Ta duy trì mỗi đường nặng bằng một cây đoạn hoặc một cây nhị phân ảo.
Ta mượn tên Auxiliary Tree để gọi mỗi cây đoạn hoặc cây nhị phân ảo này.

Với thao tác sửa trọng số cạnh, chỉ cần sửa \texttt{cost} của đỉnh tương ứng và
duy trì Auxiliary Tree chứa nó.

Với truy vấn trọng số cạnh lớn nhất trên đường giữa $u$ và $v$, ta dùng cách
ngây thơ để tìm tổ tiên chung gần nhất của hai Auxiliary Tree tương ứng với
$u$ và $v$. Theo Bổ đề 3.1, từ bất kỳ đỉnh nào lên gốc đều đi qua nhiều nhất
$\lg n$ đường nặng, nên nhiều nhất $O(\log n)$ Auxiliary Tree. Trong quá trình
tìm tổ tiên chung gần nhất, trên mỗi Auxiliary Tree đi qua, ta có thể lấy
trọng số cạnh lớn nhất trên phần đường đi thuộc Auxiliary Tree đó trong
$O(\log n)$. Cuối cùng tính thêm phần đường đi nằm trong Auxiliary Tree chứa
tổ tiên chung gần nhất. Giá trị lớn nhất trong tất cả các giá trị này là đáp
án.

Vì sửa trọng số mất $O(\log n)$, còn trả lời truy vấn mất $O(\log^2 n)$, tổng
độ phức tạp là $O(n+q\log^2 n)$.

\subsection{Dùng Splay Tree để duy trì đường đi}

Trong lời giải này, chỉ cần thay cấu trúc dữ liệu duy trì đường đi ở lời giải
trước bằng Splay Tree. Dùng cùng hàm thế năng như trong phân tích độ phức tạp
của cây động, có thể chứng minh thời gian khấu hao trả lời truy vấn là
$O(\log n)$. Vì thế năng ban đầu là $O(n\log n)$, tổng độ phức tạp là
$O((n+q)\log n)$.

Nhưng trong Bảng~\ref{tab:runtime}, ta thấy thời gian chạy của chương trình này
không lý tưởng. Điều đó dẫn tới câu hỏi sau.

\subsection{Tìm kiếm một cấu trúc dữ liệu duy trì đường đi tốt hơn}

Có tồn tại một cấu trúc dữ liệu để duy trì đường đi, không cần mạnh như
Splay Tree, nhưng đồng thời có độ phức tạp lý thuyết như lời giải 3 và hiệu
quả thực tế tốt như lời giải 2 hay không?

Trong phần thảo luận trên, ta thấy dùng cây đoạn hoặc cây nhị phân ảo để duy
trì đường đi thì mỗi Auxiliary Tree đều cân bằng tuyệt đối, nhưng độ phức tạp
lý thuyết lại cao hơn dùng Splay Tree. Vì sao?

Từ phân tích độ phức tạp khi dùng Splay Tree duy trì Auxiliary Tree, có thể
thấy Splay Tree không xem từng Auxiliary Tree một cách cô lập, mà thông qua
Path Parent xem các Auxiliary Tree như một tổng thể. Điều này đem lại cân bằng
toàn cục, nên mỗi truy vấn có độ phức tạp khấu hao $O(\log n)$.

Cây đoạn hoặc cây nhị phân ảo không tận dụng liên hệ giữa các Auxiliary Tree.
Dù từng Auxiliary Tree cân bằng đến đâu, việc mất phương hướng toàn cục lại
dẫn tới mất cân bằng trên tổng thể.

Vì vậy, nếu tận dụng Path Parent giữa các Auxiliary Tree và tùy chỉnh hình dạng
phù hợp cho mỗi Auxiliary Tree nhằm đạt cân bằng toàn cục, ta có hy vọng thu
được thuật toán đang tìm kiếm.

\subsection{Lời giải dùng cây nhị phân cân bằng toàn cục}

Theo suy nghĩ ở phần trước, ta xem tất cả Auxiliary Tree như một tổng thể và
cần tìm cách chọn hình dạng phù hợp cho mỗi Auxiliary Tree. Vì phương pháp này
đạt cân bằng trên tổng thể, ta gọi nó là \term{cây nhị phân cân bằng toàn cục}.

\subsubsection{Cây ảo}

Trước hết, ta xét từ góc nhìn toàn cục, chưa quan tâm hình dạng cụ thể của mỗi
Auxiliary Tree mà chỉ xét quan hệ giữa chúng.

Ta gọi con của một đỉnh trong Auxiliary Tree là \term{con thật} của nó. Nếu gốc
của một Auxiliary Tree là $r$ và
$\operatorname{path\mbox{-}parent}[r]=u$, ta gọi $r$ là một \term{con ảo} của
$u$. Xét một cây có gốc là gốc của Auxiliary Tree cao nhất, được nối bởi cả
cạnh thật và cạnh ảo; nhắc lại cách nhìn trong lời giải 2 và Link-Cut Trees,
ta gọi cây này là \term{cây ảo}. Cây con của một đỉnh được định nghĩa là cây
gồm chính nó và các cây con nối với nó bằng cạnh thật hoặc cạnh ảo; con trái
và con phải là con của nó trong Auxiliary Tree chứa nó; con ảo là con nối với
nó bằng cạnh ảo.

Bóng dáng của cây ảo xuất hiện trong mọi phương pháp đã thảo luận. Trong lời
giải 1 và 3, ta duy trì một cây ảo động. Trong lời giải 2, ta duy trì một cây
ảo tĩnh, nhưng độ sâu của nó là $O(\log^2 n)$.

Trong lời giải 4, mục tiêu là tìm một cây ảo có độ sâu $O(\log n)$. Tất nhiên
có thể có nhiều cây như vậy; phần 4.7.4 sẽ thảo luận thêm. Ở đây, ta chỉ gọi
những cây ảo thỏa định nghĩa cây nhị phân cân bằng toàn cục dưới đây là cây
nhị phân cân bằng toàn cục.

\subsubsection{Định nghĩa và xây dựng cây nhị phân cân bằng toàn cục}

Làm thế nào để xây dựng một cây nhị phân cân bằng toàn cục? Trước hết giả sử
ta đã dùng phân rã nặng-nhẹ để thu được một cây ảo, và đã xây dựng thành công
một cây nhị phân cân bằng toàn cục. Ta nghiên cứu các tính chất của cây này.

\paragraph{Yêu cầu 4.1.}
Trong cây nhị phân cân bằng toàn cục, với mọi cây con, tổng số đỉnh trong cây
con trái và cây con phải của nó đều không vượt quá một nửa tổng số đỉnh của
cây con đó.

Yêu cầu này thực ra rất dễ thỏa. Giả sử các đỉnh trong Auxiliary Tree chứa gốc
của cây con hiện tại, sau khi sắp theo khóa tăng dần, là
$v_1,v_2,\ldots,v_k$, trong đó $k$ là số đỉnh của Auxiliary Tree đó. Đặt
\[
  w(u)=1+\text{tổng số đỉnh trong mọi cây con ảo có cha là }u,
\]
và
\[
  s_i=\sum_{j=1}^{i} w(v_j).
\]
Khi đó $s_k$ là tổng số đỉnh của cây con này. Vì $s_0=0$ và $s_{k-1}<s_k$, tồn
tại một số nguyên $t$ sao cho
\[
  s_{t-1} < s_k/2 \le s_t.
\]
Chọn $v_t$ làm gốc của Auxiliary Tree này. Khi đó tổng số đỉnh của cây con
trái là $s_{t-1}<s_k/2$, còn tổng số đỉnh của cây con phải là
$s_k-s_t\le s_k/2$.

\paragraph{Tính chất 4.2.}
Trong cây nhị phân cân bằng toàn cục, với mọi cây con, tổng số đỉnh trong mỗi
cây con ảo của nó nhỏ hơn một nửa tổng số đỉnh của cây con đó.

\paragraph{Chứng minh.}
Theo Tính chất 3.1, tổng số đỉnh của mọi cây con ảo nhỏ hơn một nửa hiệu giữa
tổng số đỉnh của cây con hiện tại và tổng số đỉnh của cây con trái của nó, do
đó cũng nhỏ hơn một nửa tổng số đỉnh của cây con hiện tại.

\paragraph{Yêu cầu 4.3 (cấu trúc con).}
Mọi cây con của một cây nhị phân cân bằng toàn cục cũng là cây nhị phân cân
bằng toàn cục.

\paragraph{Định lý 4.1 (định lý cân bằng toàn cục).}
Độ sâu lớn nhất có thể của một cây nhị phân cân bằng toàn cục có $n$ đỉnh, nếu
gốc có độ sâu $0$, thỏa
\[
  \operatorname{max\mbox{-}depth}(n)\le \lg n.
\]

\paragraph{Chứng minh.}
Theo Yêu cầu 4.1, Tính chất 4.2 và Yêu cầu 4.3, số đỉnh trong cây con trái,
cây con phải hoặc cây con ảo của gốc đều nhiều nhất bằng một nửa ban đầu, và
mỗi cây con đó vẫn là một cây nhị phân cân bằng toàn cục. Do đó
\[
  \operatorname{max\mbox{-}depth}(n)
  \le \operatorname{max\mbox{-}depth}(n/2)+1
  \le \lg n.
\]

Từ các nhận xét trên, ta có thể viết quy trình xây dựng cây nhị phân cân bằng
toàn cục.

\begingroup
\small
\begin{verbatim}
function Build(l, r)
    tim nhi phan t sao cho
        s[t-1] < (s[r] + s[l-1]) / 2 <= s[t]
    if t > l then
        left-child[t] <- Build(l, t - 1)
    end if
    if t < r then
        right-child[t] <- Build(t + 1, r)
    end if
    return t
end function

procedure Construct-Auxiliary-Tree(v[1..k])
    tinh s[0..k]
    Build(1, k)
end procedure

procedure Compute-Weight(node)
    size(node) <- 1
    for all canh (node, v) do
        Compute-Weight(v)
        size(node) <- size(node) + size(v)
    end for
    w(node) <- size(node)
    if ton tai canh nang (node, v) then
        w(node) <- w(node) - size(v)
    end if
end procedure

procedure Make-Global-Balanced-BST(tree)
    Heavy-Light-Decomposition(tree)
    Compute-Weight(root)
    for all duong nang (v[1..k]),
        trong do v[i] la cha cua v[i+1] do
        Construct-Auxiliary-Tree(v[1..k])
    end for
end procedure
\end{verbatim}
\endgroup

Từ mã giả có thể thấy quá trình xây dựng cây nhị phân cân bằng toàn cục mất
$O(n\log n)$. Hình 2 trong bản gốc minh họa một cây, phân rã nặng-nhẹ của nó,
và cây nhị phân cân bằng toàn cục tương ứng; bản dịch không vẽ lại vì trích
xuất PDF làm hỏng bố cục hình.

\subsubsection{Lời giải dùng cây nhị phân cân bằng toàn cục}

Với thao tác sửa trọng số mà đề bài yêu cầu, ta chỉ cần duy trì Auxiliary Tree
tương ứng. Quá trình này không nhắc lại nữa. Vì độ sâu của cây nhị phân cân
bằng toàn cục là $O(\log n)$, thao tác này có độ phức tạp $O(\log n)$.

Với truy vấn trọng số cạnh lớn nhất trên đường đi, hãy nhớ lại cách làm trong
lời giải 2. Ta chia truy vấn thành các phần trên những Auxiliary Tree ``đi qua''
và phần trên Auxiliary Tree chứa tổ tiên chung gần nhất. Phần truy vấn trong
Auxiliary Tree chứa tổ tiên chung gần nhất có thể trả lời trong $O(\log n)$,
nên nó không ảnh hưởng đến độ phức tạp của lời giải. Ta tập trung vào xử lý
các Auxiliary Tree đi qua.

Trên mỗi Auxiliary Tree đi qua, ta chỉ cần biết giá trị \texttt{cost} lớn nhất
từ đỉnh hiện tại, qua cha của nó, \ldots, đến đỉnh cao nhất của đường nặng.
Trong Auxiliary Tree, điều này tương đương lấy \texttt{cost} lớn nhất trong
các đỉnh có khóa không vượt quá khóa của đỉnh hiện tại. Rõ ràng, chỉ cần bắt
đầu từ gốc của Auxiliary Tree đi xuống tới đỉnh hiện tại và dùng các giá trị
\texttt{maxcost} đã duy trì là tính được giá trị cần tìm. Nếu độ sâu của gốc
Auxiliary Tree là $\operatorname{depth}(\operatorname{root})$, còn độ sâu của
đỉnh hiện tại là $\operatorname{depth}(\operatorname{node})$, truy vấn này tốn
\[
  O(\operatorname{depth}(\operatorname{root})
    -\operatorname{depth}(\operatorname{node})+1).
\]

Giả sử khi xử lý đỉnh truy vấn $v_0$, các Auxiliary Tree đi qua lần lượt có
gốc $r_0,r_1,\ldots,r_k$, và Path Parent của $r_i$ là $v_{i+1}$. Rõ ràng ta
cần lần lượt tìm \texttt{cost} lớn nhất trên đoạn từ $v_i$ đến $r_i$.

Theo kết luận trên, tổng thời gian trên các Auxiliary Tree đi qua là
\[
\begin{aligned}
O\left(\sum_{i=0}^{k}
(\operatorname{depth}(r_i)-\operatorname{depth}(v_i)+1)\right)
&=
O\left(\sum_{i=0}^{k}
(\operatorname{depth}(v_{i+1})-\operatorname{depth}(v_i))\right) \\
&= O(\operatorname{depth}(v_{k+1})-\operatorname{depth}(v_0)) \\
&= O(\log n).
\end{aligned}
\]
Vì vậy, một truy vấn được trả lời trong $O(\log n)$.

Cũng có thể hiểu trực quan quá trình trả lời truy vấn như sau: ta tìm tổ tiên
chung gần nhất của hai đỉnh trên cây nhị phân cân bằng toàn cục này (tổ tiên
chung đó không nhất thiết là tổ tiên chung gần nhất của hai đỉnh trên cây gốc),
và thu thập thông tin liên quan trên các đỉnh đi qua. Vì độ sâu của cây nhị
phân cân bằng toàn cục là $O(\log n)$, tổng thời gian trả lời truy vấn là
$O(\log n)$.

\subsubsection{Cây nhị phân ``cân bằng toàn cục''}

Trong cây nhị phân cân bằng toàn cục, điều kiện then chốt ảnh hưởng đến độ sâu
thực ra là Tính chất 4.2. Điều kiện này có thể nới lỏng thành: với mọi cây con,
tổng số đỉnh trong mỗi cây con ảo nhỏ hơn $c$ lần tổng số đỉnh của cây con đó,
trong đó $c<1$ là một hằng số. Ta gọi cây thỏa điều kiện này là một cây nhị
phân ``cân bằng toàn cục''. Rõ ràng Yêu cầu 4.1 và 4.3 vẫn có thể thỏa mãn.
Tương tự Định lý 4.1, ta chứng minh được độ sâu của cây nhị phân ``cân bằng
toàn cục'' không vượt quá
\[
  \log_{1/c} n = O(\log n),
\]
chỉ là khi $c$ lớn thì hằng số trong độ sâu cũng lớn.

Một cây nhị phân ``cân bằng toàn cục'' hiển nhiên cũng có thể dùng để giải bài
này. Nhưng trên thực tế, ta có thể bảo đảm rằng hằng số $c$ của cây nhị phân
``cân bằng toàn cục'' xây dựng được là ít nhất $1/2$. Vì vậy phương pháp trên
không còn chỗ cải thiện. Ngay cả nếu tìm mọi cách để đạt một $c<1/2$, cũng vô
ích: khi đó yếu tố khống chế độ sâu đã chuyển thành Yêu cầu 4.1, và xét trường
hợp toàn bộ cây là một đường thẳng thì yêu cầu này không thể cải thiện nữa.
Do đó, cây nhị phân cân bằng toàn cục đã trình bày ở trên là đủ tốt.

\subsection{So sánh và tổng kết}

Phân tích Bảng~\ref{tab:runtime}, ta thấy lời giải 3 và lời giải 2 tuy xuất
phát từ cùng một hướng, nhưng do tính đặc biệt trong phân tích độ phức tạp của
Splay Tree, lời giải 3 có độ phức tạp lý thuyết thấp hơn. Chỉ vì hằng số của
Splay Tree khá lớn nên trong chạy thực tế, chương trình của lời giải 3 chậm
hơn nhiều so với chương trình của lời giải 2.

Đồng thời, thời gian chạy của lời giải 1 và lời giải 3 rất gần nhau. Điều này
cho thấy thao tác đổi Preferred Child trong Link-Cut Trees thực tế rất nhanh.
Phần tốn thời gian nhất của Link-Cut Trees chính là thao tác dùng Splay Tree để
duy trì đường đi. Điều thú vị là sau khi dùng Splay Tree làm cấu trúc dữ liệu
duy trì đường đi, sự xuất hiện và ứng dụng của Link-Cut Trees trở nên rất tự
nhiên: nó chỉ trả thêm một chi phí rất nhỏ nhưng nhận được khả năng mở rộng
rất lớn. Khả năng mở rộng này chính là do cấu trúc dữ liệu đặc biệt Splay Tree
và cách phân tích độ phức tạp đặc biệt của nó đem lại. Cùng với khả năng mở
rộng lớn đó, Splay Tree cũng không tránh khỏi hằng số đáng kể. Có thể tạm xem
việc dùng Splay Tree duy trì đường đi là thiết kế dành riêng cho Link-Cut Trees;
nếu không thì có phần dùng quá mức.

Điều này gợi ý ta suy nghĩ xem có tồn tại một cấu trúc dữ liệu khác cho trường
hợp hình dạng cây không thay đổi, dùng một cây cân bằng không phải trả chi phí
xoay cao để duy trì đường đi, mà vẫn có độ phức tạp lý thuyết xấu nhất giống
lời giải 1 và 3 hay không. Suy nghĩ theo hướng đó dẫn tới lời giải 4 ở đây.

Từ lời giải 3 và 4 có thể thấy, khi giải một số bài toán, tối ưu cục bộ chưa
chắc đem lại tối ưu toàn cục. Ta không nên cố chấp theo đuổi tối ưu cục bộ, mà
cần luôn giữ góc nhìn toàn cục, nắm được liên hệ giữa tổng thể và cục bộ, để
từ đó đạt tối ưu trên tổng thể. Đây chính là mục đích chính của tác giả khi
viết bài này.

\appendix
\section{Lời cảm ơn}

Bạn Yuan Xinhao của Trường Trung học Changjun, Hồ Nam, đã giúp tác giả rất
nhiều trong quá trình hoàn thành báo cáo lời giải này. Xin gửi lời cảm ơn tới
bạn ấy.

\begin{thebibliography}{9}
\bibitem{st-dt}
Daniel D. Sleator, Robert Endre Tarjan,
\textit{A data structure for dynamic trees},
Journal of Computer and System Sciences, vol.~26, no.~3, pp.~362--391,
June 1983.

\bibitem{splay}
Daniel Dominic Sleator, Robert Endre Tarjan,
\textit{Self-adjusting binary trees},
Proceedings of the fifteenth annual ACM Symposium on Theory of Computing,
pp.~235--245, December 1983.

\bibitem{zhou}
Zhou Gelin,
\textit{Báo cáo lời giải ``Query on a tree (SPOJ375)''}.
\end{thebibliography}

\end{document}
